\documentclass{article}

\usepackage{amsmath,amssymb}
\usepackage{xurl}
\usepackage[hidelinks]{hyperref}

\input{command}

\title{Wolf Theorem~6.16: The Two Schwarz Orientations in the Printed Proof}
\author{}
\date{2026-08-24}

\begin{document}
\maketitle
\gapnote{local-correction}{resolved}

\section{Statement and invoked fixed-point theorem}

Let $T:\MN{d}\to\MN{d}$ be positive and trace preserving, and let $T^*$ be its
trace-pairing adjoint. Theorem~6.16 of Ref.~\cite{Wolf2012Quantum} assumes that
$T$ satisfies the Schwarz inequality. It concludes that the peripheral space
has the weighted matrix-block form in source Equations~(6.66)--(6.67), and
that $T$ acts through the block permutation and unitary conjugations in source
Equation~(6.68). The local source is
\path{Notes/WolfNoteTexSource/ch06_spectral_properties.tex},
lines~1592--1624.

The proof takes a subsequential limit of powers,
\begin{align}
  I
    &=\lim_{i\to\infty}T^{n_i},
  \label{eq:wolf_theorem6_16_schwarz_orientation_recurrent_projection}
\end{align}
and states that $I$ remains Schwarz because the Schwarz property is preserved
under composition. It later excludes matrix transposition because $T$, and
therefore each induced block map $T_k$, is Schwarz. These steps use the
Schwarz inequality for $T$ itself; see local source lines~1626--1633 and
1660--1663.

The same proof identifies the peripheral space with $\operatorname{Fix}(I)$
and invokes Theorem~6.14 of Ref.~\cite{Wolf2012Quantum} to obtain the weighted
matrix-block structure. That theorem assumes that the trace adjoint of its
trace-preserving map is Schwarz. Applied to $I$, the required hypothesis is
\begin{align}
  I^*(A^\dagger A)
    &\succeq I^*(A^\dagger)I^*(A)
    \qquad\text{for every }A\in\MN{d}.
  \label{eq:wolf_theorem6_16_schwarz_orientation_required_adjoint_schwarz}
\end{align}
The local source for this hypothesis is lines~1467--1474. Closure under powers
gives the Schwarz inequality for $I$ from the corresponding inequality for
$T$, but it does not give
(\ref{eq:wolf_theorem6_16_schwarz_orientation_required_adjoint_schwarz}).
The printed proof supplies no argument that changes the orientation.

\section{Corrected proof contract}

The minimal explicit correction that validates the printed proof is to assume
both Schwarz inequalities:
\begin{align}
  T(A^\dagger A)
    &\succeq T(A^\dagger)T(A)
    \qquad\text{for every }A\in\MN{d},
  \label{eq:wolf_theorem6_16_schwarz_orientation_map_schwarz}\\
  T^*(A^\dagger A)
    &\succeq T^*(A^\dagger)T^*(A)
    \qquad\text{for every }A\in\MN{d}.
  \label{eq:wolf_theorem6_16_schwarz_orientation_adjoint_schwarz}
\end{align}
The first orientation is used for the block dynamics and the exclusion of
transposition. The second is stable under powers and finite-dimensional
limits. Indeed,
\begin{align}
  (T^n)^*
    &=(T^*)^n,
  \label{eq:wolf_theorem6_16_schwarz_orientation_adjoint_power}\\
  I^*
    &=\lim_{i\to\infty}(T^*)^{n_i}.
  \label{eq:wolf_theorem6_16_schwarz_orientation_adjoint_limit}
\end{align}
Thus
(\ref{eq:wolf_theorem6_16_schwarz_orientation_adjoint_schwarz}) implies
(\ref{eq:wolf_theorem6_16_schwarz_orientation_required_adjoint_schwarz}), and
Theorem~6.14 may be applied to $I$ with its literal source hypothesis.

This correction does not claim that the conclusion of Theorem~6.16 is false
under the single printed Schwarz hypothesis. A different proof might establish
the required fixed-point structure. The correction records only the contract
supported by the proof as written. In particular, a formal theorem following
that route cannot infer that the trace-pairing adjoint of an arbitrary Schwarz
map is Schwarz.

\section{Formalization boundary}

The source-general fixed-point theorem has the literal interface of
Theorem~6.14 of Ref.~\cite{Wolf2012Quantum}: the map $S$ is positive and trace
preserving, and $S^*$ satisfies the Schwarz inequality.\footnote{Formal
declaration:
\leanid{IsPositiveMap.exists_fixedPoints_densityBlocks_with_zero} in
\doclink{QICLean/Channel/FixedPoint/WolfTheorem614}.}

The completed source-facing assembly for Theorem~6.16 exposes both
orientations. Its recurrent-projection component passes the adjoint orientation
to the fixed-point theorem, while its block-classification component retains
the orientation for the original map.\footnote{Formal declaration:
\leanid{IsPositiveMap.exists_wolfTheorem616} in
\doclink{QICLean/Channel/Peripheral/StructureOfCycles}. It preserves the exact
zero-extended density-block membership form and proves the output-indexed
unitary action in source Equation~(6.68) of
Ref.~\cite{Wolf2012Quantum}.}
The construction is organized by \ghissue{40}; this contract audit is
\ghissue{407}.

\section{Conclusion}

The printed proof of Theorem~6.16 uses both Schwarz orientations but states
only the one for $T$ \cite[Theorem~6.16]{Wolf2012Quantum}. Adding the Schwarz
inequality for $T^*$ is a local hypothesis correction that validates the
invocation of Theorem~6.14 without changing the source terminology, block
notation, or proof order. Unless a separately sourced derivation removes the
extra assumption, both orientations remain the source-facing formal contract.

\bibliographystyle{alpha}
\bibliography{references}

\end{document}
